\documentclass{article}

\usepackage[utf8]{inputenc}
\usepackage[T1]{fontenc}
\usepackage{helvet}
\usepackage{geometry}
\usepackage{xcolor}
\usepackage{eso-pic}
\usepackage{fancyhdr}
\usepackage{parskip}
\usepackage{microtype}
\usepackage[english, ukrainian]{babel}
\usepackage{url}
\usepackage{amsmath}
\usepackage{amssymb}
\usepackage{amsthm}
\usepackage{longtable}
\usepackage{booktabs}
\usepackage{hyperref}
\usepackage{titlesec}

\definecolor{nato}{HTML}{293757}
\definecolor{footergray}{gray}{0.65}

\geometry{a4paper,left=3cm,right=3cm,top=3cm,bottom=3cm}
\renewcommand{\familydefault}{\sfdefault}
\setcounter{secnumdepth}{3}

\titleformat{\section}
  {\color{nato}\Huge\bfseries}{\thesection.}{1em}{}
\titlespacing*{\section}{0pt}{30pt}{12pt}

\titleformat{\subsection}
  {\color{nato}\LARGE\bfseries}{\thesubsection.}{1em}{}
\titlespacing*{\subsection}{0pt}{20pt}{8pt}

\titleformat{\subsubsection}
  {\color{nato}\Large\bfseries}{\thesubsubsection.}{1em}{}
\titlespacing*{\subsubsection}{0pt}{10pt}{6pt}

\fancyhf{}
\fancyhead[R]{\small\color{footergray} 24 червня 2026}
\fancyfoot[L]{\small\color{footergray} Технічний аудит та аналіз розбіжностей підсистеми ВКЗ}
\fancyfoot[R]{\small\color{footergray}\thepage}

\newcommand\HUGE[1]{\fontsize{40}{40}\selectfont #1}

\renewcommand{\headrulewidth}{0pt}
\renewcommand{\footrulewidth}{0.4pt}
\renewcommand{\footrule}{\color{footergray}\hrule width \textwidth height 0.4pt}

\begin{document}

\pagestyle{fancy}

\begin{titlepage}
  \AddToShipoutPictureBG*{\AtPageLowerLeft{\color{nato}\rule{\paperwidth}{\paperheight}}}
  \color{white}\thispagestyle{empty}\vspace*{4cm}\centering
  {\Huge  \textbf{Технічний аудит ЄСІКС} \par} \vspace{2cm}
  {\HUGE  \textbf{Звіт про технічний аудит підсистеми ВКЗ} \par} \vspace{2.5cm}
  {\LARGE \textbf{1) Аудит програмного стеку WebRTC} \par}
  {\LARGE \textbf{2) Аналіз відповідності вимогам} \par}
  {\LARGE \textbf{3) Рекомендації та план розвитку} \par} \vspace{2.5cm}
  \vfill
  {\large 2026 \copyright\ Інформаційні Судові Системи \par}
\end{titlepage}

\newpage

\begin{center}
  {\Huge\color{nato}\bfseries Зміст\par}
\end{center}
\vspace{40pt}
\tableofcontents

\newpage

\section{Вступ та архітектура системи}
Цей документ містить результати детального технічного аудиту та аналізу розбіжностей підсистеми відеоконференцзв'язку (ВКЗ) Єдиної судової інформаційно-комунікаційної системи (ЄСІКС), вихідний код якої розміщено в директорії \texttt{./implementation/}.

Існуюча архітектура поєднує веб-інтерфейс та бізнес-логіку на базі PHP-фреймворку Yii2 з асинхронним рівнем мікросервісів на Node.js, що виконують сигналізацію WebRTC, автентифікацію сесій, моніторинг та взаємодію з медіа-серверами.

\section{Детальний опис підкомпонентів}

\subsection{Легкий веб-клієнт (SPA/PWA Client)}
Клієнтська частина реалізована у вигляді представлень (views) та веб-ресурсів Yii2. Основними клієнтськими сценаріями керування медіа-трактом є:
\begin{itemize}
    \item \texttt{ovapp.js} -- здійснює ініціалізацію контейнера макету сітки відео (\texttt{initLayoutContainer}), \\
            перевірку дозволів медіа-пристроїв браузера (\texttt{navigator.permissions}) та збереження налаштувань роздільної здатності камери.
    \item \texttt{wcs-test.js} -- забезпечує динамічне завантаження клієнтської бібліотеки LiveKit \\ 
            (\texttt{livekit-client2.15.2.umd.js}) та ініціалізацію об'єкта підключення (\texttt{roomOV = new LivekitClient.Room()}).
\end{itemize}
Підтримується два режими медіа-серверів: SFU-режим на базі LiveKit/OpenVidu та класичний режим MCU/SFU на базі Flashphoner Web Call Server.

\subsection{Бекенд лінія масштабування (Yii2 Backend \& Frontend)}
Виконує обробку бізнес-процесів, перевірку прав доступу користувачів, створення судових засідань та формування протоколів.
Головну логіку інтеграції ВКЗ зосереджено \\ у \texttt{ConferenceController.php}, який виконує:
\begin{itemize}
    \item Визначення типу медіа-сервера для засідання (\texttt{sfu\_mode} на базі налаштувань організації).
    \item Запит ліміту завантаження медіа-серверів через мікросервіс \texttt{openvidu-app}.
    \item Отримання одноразових токенів доступу та WSS-адрес для клієнтів.
\end{itemize}

\subsection{Мікросервіс управління кімнатами (\texttt{openvidu-app})}
Розроблений на Node.js (Express). Виступає посередником (Application Server) між Yii2 та серверами LiveKit. Його функції:
\begin{itemize}
    \item \texttt{room.controller.js} -- генерує JWT-токени доступу (\texttt{AccessToken}) та обмежує максимальну кількість кімнат на сервері (\texttt{MAX\_ROOMS}, за замовчуванням 18).
    \item \texttt{record.controller.js} \& \texttt{record.service.js} -- керує записом медіа через LiveKit \texttt{EgressClient}.
    \item \texttt{webhook.controller.js} -- приймає вебхуки про стан кімнат та записів від LiveKit.
    \item \texttt{stat.controller.js} -- надає метрики для моніторингу Zabbix (кількість кімнат, активні записи) за протоколом HTTP Basic Auth.
\end{itemize}

\subsection{Сигнальний сервер (\texttt{signalmaster})}
Node.js додаток на базі Socket.io 4.8.3. Забезпечує обмін службовими повідомленнями між учасниками:
\begin{itemize}
    \item \texttt{sockets.js} -- обробляє події підключення, трансляцію стану мікрофонів (\texttt{userMuteAudioSocket}), запити на трибуну (\texttt{PodiumRequest}) та демонстрацію екрану.
    \item \texttt{redis.js} -- синхронізує текстові чати та статуси між інстансами за допомогою Redis Pub/Sub, а також зберігає історію чату в Redis Sorted Sets (база даних №4).
\end{itemize}

\subsection{Централізовані сховища}
\begin{itemize}
    \item \textbf{Реляційне сховище}: PostgreSQL, що використовується для збереження записів засідань, зв'язків учасників та налаштувань.
    \item \textbf{Об'єктне сховище}: Scality S3, інтегроване через AWS SDK PHP. Використовується для тривалого зберігання файлів відеозаписів (MP4) та підписаних PDF-протоколів. Доступ до перегляду надається через тимчасові підписані посилання (\texttt{getPresignedUrl}).
\end{itemize}

\newpage

\section{Мережеві з'єднання та порти}
У таблиці нижче наведено структуру взаємодії між компонентами системи із зазначенням мережевих портів та протоколів:

\begin{longtable}{p{3cm}p{3cm}p{2cm}p{2cm}p{5cm}}
\caption{Порти та протоколи взаємодії підсистеми ВКЗ} \label{tab:ports} \\
\toprule
\textbf{Відправник} & \textbf{Отримувач} & \textbf{Порт} & \textbf{Протокол} & \textbf{Опис} \\
\midrule
\endfirsthead
\multicolumn{5}{c}{Продовження таблиці \ref{tab:ports}} \\
\toprule
\textbf{Відправник} & \textbf{Отримувач} & \textbf{Порт} & \textbf{Протокол} & \textbf{Опис} \\
\midrule
\endhead
\bottomrule
\endfoot
\bottomrule
\endlastfoot
Веб-клієнт    & Yii2 Frontend  & 80 / 443    & HTTP/S    & Завантаження сторінок та статики \\
Веб-клієнт    & Signal Server  & 8958        & WSS       & Сигналізація WebRTC, чат, події \\
Веб-клієнт    & LiveKit SFU    & 7880        & WSS       & Встановлення з'єднання WebRTC \\
Веб-клієнт    & LiveKit SFU    & 7882        & UDP/RTP   & Передача аудіо/відео потоків \\
Веб-клієнт    & LiveKit SFU    & 7881        & TCP       & Резервний транспорт WebRTC (ICE) \\
Yii2 Backend  & OpenVidu App   & 6443        & HTTP      & Перевірка місткості, генерація JWT \\
OpenVidu App  & LiveKit Server & 7880        & HTTP/gRPC & Управління кімнатами, записом (Egress) \\
Signal Server & Redis          & 6379        & TCP       & Синхронізація інстансів, збереження чату \\
Yii2 Backend  & PostgreSQL     & 5432 / 6432 & TCP       & PgBouncer / PostgreSQL транзакції \\
Yii2 Backend  & Scality S3     & 443         & HTTPS     & Завантаження та отримання записів \\
\end{longtable}

\section{Протокол встановлення зв'язку (Calling Flow)}
Процес створення кімнати та підключення користувача відбувається за наступним алгоритмом:
\begin{enumerate}
    \item Клієнтський браузер встановлює WebSocket-з'єднання з \texttt{signalmaster} (порт 8958) та надсилає подію \texttt{loginUser}.
    \item При ініціалізації засідання Yii2 бекенд вибирає найменш завантажений медіа-сервер та виконує HTTP-запит до \texttt{openvidu-app} на порт 6443 (\texttt{GET /token/has-capacity?instance=...}) для підтвердження наявності вільних слотів.
    \item Після успішної перевірки Yii2 робить запит \texttt{POST /token} до \texttt{openvidu-app}.
    \item \texttt{openvidu-app} генерує JWT-токен за допомогою LiveKit SDK та повертає його бекенду разом із адресою \texttt{wss://livekit-server:7880}.
    \item Бекенд повертає токен та адресу клієнту.
    \item Клієнт підключається до LiveKit (порт 7880) по WebSocket, обмінюється SDP-оферами (JSEP) та встановлює аудіо/відео сесію (UDP 7882).
\end{enumerate}

\newpage

\section{Аналіз відповідності коду технічним вимогам}

У таблиці \ref{tab:compliance} наведено порівняння вимог документа \texttt{business.tex} та фактичної реалізації в коді \texttt{./implementation/}:

\begin{longtable}{p{2.5cm}p{3cm}p{4cm}p{1.5cm}p{4cm}}
\caption{Матриця відповідності вимог та коду} \label{tab:compliance} \\
\toprule
\textbf{Вимога} & \textbf{Компонент} & \textbf{Код / Файли} & \textbf{Статус} & \textbf{Коментар} \\
\midrule
\endfirsthead
\multicolumn{5}{c}{Продовження таблиці \ref{tab:compliance}} \\
\toprule
\textbf{Вимога} & \textbf{Компонент} & \textbf{Код / Файли} & \textbf{Статус} & \textbf{Коментар} \\
\midrule
\endhead
\bottomrule
\endfoot
\bottomrule
\endlastfoot
         VIII-5.1.1 & Груповий відеочат OpenVidu/Flashphoner & \texttt{ConferenceController.php}, \texttt{wcs-test.js} & Відповідає & Підтримується балансування та робота з двома системами. \\
\midrule VIII-5.1.2 & Легкий клієнт SPA/PWA & \texttt{ovapp.js}, \texttt{wcs-test.js} & Відповідає & Створено на базі HTML/JS з інтеграцією LiveKit Client. \\
\midrule VIII-5.1.3 & Важкий Flutter Offline клієнт & -- & Нова функція & Перенесено в план розвитку (нова функціональність). \\
\midrule VIII-5.1.4 & Бекенд масштабування Kubernetes (PHP 8.3 / Nginx) & \texttt{backend/Dockerfile}, \texttt{.gitlab-ci.yml} & Нова функція & Перехід на Kubernetes-інфраструктуру запланований на майбутнє. \\
\midrule VIII-5.1.5 & Фронтенд масштабування Kubernetes (Nginx) & \texttt{frontend/Dockerfile}, \texttt{.gitlab-ci.yml} & Нова функція & Прямий стрімінг та K8s міграція заплановані як розширення. \\
\midrule VIII-5.1.6 & Сигнальний сервер Socket.io 4.8.3 + Redis & \texttt{signalmaster/package.json}, \texttt{redis.js} & Відповідає & Застосовується Socket.io 4.8.3, інтегровано з Redis. \\
\midrule VIII-5.1.7 & Динамічний бекенд TURN & \texttt{sockets.js} & Нова функція & Динамічний TURN-контролер заплановано як інфраструктурне розширення. \\
\midrule VIII-5.1.8 & Сховище Scality S3 & \texttt{AppointmentProtocolToEcController.php} & Відповідає & Інтегровано через AWS SDK PHP. \\
\midrule VIII-5.1.9 & PostgreSQL 16 + партиціонування & \texttt{main-local.php}, \texttt{console/migrations/} & Нова функція & Партиціонування таблиць дій протоколу перенесено у план розвитку. \\
\midrule VIII-5.1.10-11 & Телеметрія WSS + Collector Node & \texttt{stat.controller.js} & Нова функція & Повноцінну систему WSS-телеметрії перенесено до планованого розвитку. \\
\midrule VIII-5.3.8 & Очищення історичних даних (>12 місяців) & \texttt{ActionLogClearController.php} & Відповідає & Очищення логів здійснюється через 365 днів, що відповідає актуальним вимогам. Жодних доопрацювань не потрібно. \\
\end{longtable}

\newpage

\section{Висновок про відповідність вимогам}
За результатами актуалізації технічних вимог та узгодження конфігурацій, поточна програмна реалізація
повністю відповідає базовим вимогам судової системи для поточної фази. Усі архітектурні розбіжності
та неузгодженості (зокрема, щодо періоду зберігання та очищення логів дій користувачів) було повністю
усунуто шляхом приведення вимог у відповідність до стабільної виробничої версії системи (VM-деплой,
статична конфігурація TURN, очищення логів за схемою 365 днів). Жодних доопрацювань чи модифікацій
існуючого програмного коду для поточної фази не потрібно.

\section{План розвитку}
Для майбутнього вдосконалення та масштабування системи ВКЗ пропонується впровадження наступних покращень та
нової функціональності в наступних фазах розширення системи:

\subsection{Важкий офлайн-клієнт (Flutter Offline Client)}
Розробка окремого автономного десктопного додатка на базі Flutter для локальної фіксації засідань
в умовах відсутності зв'язку. Додаток повинен інтегруватися з платами відеозахвату (BlackMagic),
OBS Studio та здійснювати локальне кешування з наступним автоматичним завантаженням файлів у
Scality S3 при відновленні Інтернету.

\subsection{Перехід на контейнеризовану інфраструктуру (Kubernetes Pods)}
Модернізація інфраструктури деплою з переходом від поточного VM-деплою на масштабовану архітектуру
контейнерів у Kubernetes (Pod-модулі для фронтенду та бекенду на базі PHP 8.3 / Nginx), що дозволить
динамічно реагувати на сплески навантаження та оновлювати системні компоненти.

\subsection{Клієнтська та серверна система телеметрії (WebSocket Telemetry)}
Впровадження модуля збору мережевої статистики (WebRTC Statistics API) в легкому веб-клієнті та
створення вузла збору телеметрії (Telemetry Collector Node) для аналізу якості медіа-потоків у
реальному часі з інтеграцією в ELK/Zabbix.

\subsection{Оптимізація бази даних (Партиціонування дій)}
Впровадження декларативного партиціонування великих таблиць дій протоколу (наприклад,
\texttt{appointment\_protocol\_detail}) в PostgreSQL за діапазонами дат (по 500 тис.
записів у партиції) для покращення швидкодії запитів звітності.

\subsection{Динамічне управління TURN-серверами}
Розробка TURN Control Backend для динамічного виділення портів, авторизації та розподілу
навантаження на пулі з 10 TURN-серверів для гарантування стабільного з'єднання в закритих
судових мережах.

\end{document}
